
\subsection{Preregistration}
The analysis plan was pre-registered on AsPredicted; see \url{https://aspredicted.org/5p4dq2.pdf}. 
The preregistration describes our main hypotheses, all interventions, primary and secondary 
outcomes, and statistical methods. When we conduct additional analyses beyond those described 
in the preregistration, we describe them as such.

\subsection{Trial protocol}

We collected data over three waves, indicated in Figure \ref{fig:study_flow}.
All participants were recruited via the online platform 
Prolific and participants took survesy on Qualtrics.  The study was advertised
as an academic survey asking about experiences, views, and backgrounds. We avoided
signalling our intentions or study goals. Our baseline wave measured flu vaccine 
intentions, and we collected information on health conditions predicting vulnerability to flu, 
as well as several covariates: prior vaccine experiences, trust in government vaccine 
information, whether participants follow doctor’s advice on vaccines, and the use and 
reliability of several sources for obtaining  health care information 
(such as their doctors, the CDC, newspapers, and social media). 

Our main survey wave, run the week after the conclusion of our baseline survey, was administered 
to baseline wave participants who indicated flu vaccine hesitancy,
meaning they said they “may or may not” or “do not intend” to get the flu vaccine.  In the 
main survey, we randomly assigned participants to treatment arms, 
elicited prior and posterior beliefs about side effect risk, asked about intentions to 
vaccinate, and asked whether participants found the described study trustworthy and 
(separately) relevant to their vaccination decision. We gave participants the opportunity 
to click on a link to sign up for a vaccine appointment. 
After our main survey was running for two days, we noticed that enrollment was lower
than anticipated in our sample size calculations, so we increased the payment. In our 
follow up survey (open to all main survey participants), begun a week after the 
main survey ended, we asked if participants had vaccinated against the flu in the last month. 
All survey waves included an attention check. 

\subsection{Interventions}
We randomly assigned participants to a control condition or one of three information provision 
conditions. Randomization was conducted using the randomizer tool on Qualtrics. 
All participants are first told the following:
\begin{quote}
Getting the flu vaccine is a personal decision. It carries benefits as well as risk. 
The main benefit is reducing the chance of getting severe flu. 

The main risk is an adverse event such as headache, muscle ache, tiredness, 
joint pain, fever, chills, or general discomfort, in the days after vaccination. 
“Severe” adverse events are ones that are bad enough to disrupt daily activity. 
\end{quote}
Our definition of severe adverse events corresponds roughly to a grade 3 adverse event rate.
We focus on severe events because they represent a ubstantial cost - a loss of work or leisure
- and are measured in clinical trials. 

Participants are then told that, to measure benefits and risks, vaccine manufacturers 
conduct clinical trials with a placebo arm and a vaccine arm, and in the clinical trial 
for Fluarix,  the placebo arm had a severe adverse event rate of 3 percent in the next few 
days. The control arm sees no further information. In the treatment arms, patients see 
additional information on the same screen. We test three treatment arms: 

\begin{enumerate}
  \item \textbf{Industry arm:} Participants are truthfully told that people getting the flu 
  vaccine in the trial were not statistically more likely to experience a severe adverse 
  event than people getting the placebo. 

  \item \textbf{Academic arm:} Participants are truthfully told about a study conducted 
  by “academic researchers” in which participants were offered small rewards as an 
  encouragement to vaccinate.  They are further told that the study found no significant 
  increase in severe adverse events, and the researchers concluded that the flu vaccine 
  did not increase the rate of severe adverse events. 
  (The study is \citet{sacks2025preferences}.) 

  \item \textbf{Representative arm:} The representative arm was identical to the academic 
  arm except that participants were additionally told, truthfully, that the study was 
  conducted among Prolific participants who weren’t sure they wanted to get the flu vaccine. 
\end{enumerate}
Appendix Figure \ref{fig:interventions} provides screen shots showing what partiipants see 
in each arm.

\textbf{Rationale:} Our intervention aimed to provide quantitative information on flu vaccine 
side effects from scientific studies. We expected that participants had a general understanding 
that pharmaceutical companies and researchers conduct trials to measure efficacy and side 
effects. By providing this background information to the control arm, we keep their experience 
similar  to the treatment arms, and indeed the only difference between control and industry 
arm is one additional statistic. Our industry arm reports on the Fluarix trial, because 
its vaccine arm experienced especially low adverse event rate. 

In the academic arm, we presented information from \citet{sacks2025preferences} because 
participants in that study were representative of flu vaccine hesitant prolific participants, 
and hence we could truthfully describe that study in the personal arm. Our academic and 
personal arms could have proceeded without first describing the placebo rate in the 
Fluarix trial. However presenting this information let us measure information salience 
in a consistent way across arms, as we explain below. 

Our study design decisions followed the recommendations in \citet{haaland2023designing}
and \citet{stantcheva2023run}. For example, we elicit priors as well as posteriors to facilitate 
heterogeneous effect estimation (as expected in information provision experiments).
We recruit subjects using neutral language to avoid telegraphing our intentions, and
throughout the survey we provide face-saving language, especially emphasizing
that vaccination carries costs and benefits, to reduce eperiment demand effects.
Our specific interventions were influenced by information provision experiments, 
such as \citet{chopra2025home} and especially \citet{roth2024misperceived}, 
where participants are provided with straightforward statistical information. 

\subsection{Measures}

\textbf{Primary outcome:} Our pre-specified primary outcome is the short-run personal 
side effect belief, which we conceptualize as 
the believed effect of vaccination on the short run experience of side effects severe enough 
to disrupt daily activities. To operationalize this measure, we first ask 
participants how likely they are to experience severe adverse events over 
the next few days, assuming they do not vaccinate and are not vaccinated.%
\footnote{The exact text of our question is, 
"For this question, please imagine, regardless of your actual vaccination status, 
that you \textbf{DO NOT} get the flu vaccine tomorrow or in the following few days.
Thinking about yourself, how likely do you think you would be to experience a severe 
adverse event, such as headache, muscle ache, tiredness, joint pain, fever, chills, or 
general discomfort, in the following few days? 
(\textit{As a reminder, severe means bad enough to disrupty your daily activity.})"
}
Call this
$\tilde{b}_i(0)$, the believed side effect rate absent vaccination.  Then we ask the 
same question but assuming they do vaccinate tomorrow. Call this $\tilde{b}_i(1)$, 
the believed side effect rate given vaccination.  Then the short run personal side effect 
belief is $\tilde{\Delta}_i = \tilde{y}_i(1) - \tilde{y}_i(0)$. 
This measure is person-specific, limited to a short time horizon, and a causal contrast.%
\footnote{Designing the measure as a causal contrast is consistent with the recommendations
of \citet{smith2021participants}.}

We measure personal beliefs twice, before and after information provision. To avoid asking 
the identical questions twice, our before-information-provision version uses a 
likert scale for $\tilde{b}_i(1)$ and $\tilde{b}_i(0)$. These prior beliefs are a covariate 
rather than an outcome. 

\textbf{Secondary outcomes:} We also measure beliefs about what happened in the trial.
After the information provision screen, we ask all arms (including control) 
what they think is the severe side adverse event rate in the vaccine arm of the 
manufacturer trial. We use this secondary outcome to assess the salience of the 
information provided. We incentivize the elicitation of this belief, offering participants 
an additional \$0.50 if their answer is within 1 percentage point of the truth. 
Our elicitation of personal and objective beliefs is similar to that in 
\citet{roth2024misperceived}.

We collected three additional secondary outcome metrics related to vaccination behavior. 
At the end of the main survey we ask participants if they intend to vaccinate against the flu. 
We code participants as intending to vaccinate if they answer that they are already vaccinated or 
intend to vaccinate (other possible answer: may or may not vaccinate, intend to vaccinate). 
We provide links for participants to click on to schedule a vaccination, and we track whether 
they click on the link. Third, in our follow up survey we ask participants if they have vaccinated 
in the last month.

\textbf{Trustworthiness and relevance:} After our intervention and after we elicited personal side 
effect beliefs, we asked all participants how trustworthy they found the described industry study, 
and how relevant to their decision to vaccinate they found it. For participants in the academic and 
representative arm, we also asked the analogous questions about the academic study. 
These questions, which were drawn from \citet{alsan2024representation}, were on a 0-10 scale.

\subsection{Methods}

\textbf{Sample size determination:} Our goal was to enroll 4,000 participants in the main 
survey, all of whom would be invited to complete the follow-up survey,
yielding 700 participants per arm in the follow-up trial under an assumed 30 percent attrition 
rate, for 80 percent power to detect a treatment effect of 0.15 standard deviations. 
This is the minimum power \citet{haaland2023designing} recommend for information provision 
experiments.  Simulations indicated that we would be highly powered to detect plausible effect 
sizes on beliefs (e.g. 25 percent learning rates), but we would likely be underpowered to
detect effects on vaccinations unless beliefs responded very sharply to our information
provision (e.g. learning rates above 50 percent). 
We therefore pre-specified beliefs rather than vaccinations as the primary outcome. 
To achieve our target sample size we ran the baseline survey until we had 8,000 completions. 
We prespecified that we would run the main survey for 1 week or until we had 4,000 completions. 
We ran the follow-up survey until a month had passed from the main survey ending, because our 
lookback period to measure  vaccination was one month. 

\textbf{Sample inclusion criteria:} Americans 18 and older were eligible for the baseline study.
All vaccine hesitant baseline survey participants were invited to participate in the main 
survey and all main survey participants were invited to the follow-up survey. We exclude cases 
who fail attention checks or with missing belief measures. For secondary outcome analysis 
we additionally exclude participants missing that outcome, so the sample size varies across 
secondary outcomes. We also excluded survey attempts  after the first (a possibility we did 
not anticipate and therefore did not prespecify).

\textbf{Specification:} Our pre-specified primary specification estimates the effect of each 
arm on outcomes with the following linear regression, estimated by ordinary least squares:
\begin{equation} \label{eq:main}
  y_i = \beta_0 + \beta_1 industry_i + \beta_2 academic_i + \beta_3 representative_i + 
    X_i \theta + \epsilon_i.
\end{equation}
Here $\beta_1, \beta_2$ and $\beta_3$ give the effect of industry, academic, and personal 
arm relative to control,  and $X_i$ is a vector of prespecified covariates, included to improve 
power.  Specifically, we include indicators for each level of the following categorical 
variables: prior beliefs; baseline vaccine intentions; reported vaccine experiences (separately 
for flu and COVID), with possible experiences unvaccinated, vaccinated without remembered side 
effects, vaccinated with mild side effects, or vaccinated with moderate side effects, or 
vaccinated with severe side effects); demographics (age, gender, education, income, race, 
ethnicity), political views) trust in 
government on vaccines (ranging from strongly disagree to strongly agree), and reported 
health conditions (asthma, lung disease, heart disease, diabetes, kidney disease, or 
have a condition but prefer not to say which). 
We include separate indicators for missing covariate values. 
 
